\section{A direct rank estimate on the ramified source}
\label{sec:spectral-spanning}

The short integral forms lie in a specified subspace of the source.
Each new translate of that subspace increases the allowance for denominators.
The question is how much of the target can be filled on the translate budget
of Theorem~B. Keeping the original fourth-power error estimates until the
end gives an upper bound of $(19+o(1))b^{-4}$ for the missed fraction on that budget.
The statement in \texttt{ShortAndEnough.lean} uses $4b^{-2}$.
The coefficient $19$ is optimal for the two moment inequalities and source
codimension estimate used here.

\paragraph{Inputs.}
Let $S,H$ be finite-dimensional complex Hilbert spaces, $h=\dim S>0$,
and let $G$ act unitarily on $H$ by $U_s$. An orthogonal projection $E$
commutes with every $U_s$. Put $V=EH$, $d=\dim V>0$, and let $P:S\to H$
denote the scaled Poincar\'e map. In the source coordinates let $S_r\subset S$
be the space of ramified forms and put $h_r=\dim S_r$. Write
$\mathbb E_s$ for the uniform average over $G$ and $\Pi_L$ for orthogonal
projection onto $L$. The three inputs are
\begin{align}
 \|EP\Pi_L\|_{\rm HS}^2&\ge\dim L-\alpha h
      &&\text{for every subspace }L\subset S,\label{sp:near}\\
 d\,\mathbb E_s\|P^*U_sP\|_{\rm HS}^2&\le(1+\beta)h^2,\label{sp:mom}\\
 h_r&\ge(1-\gamma)h.\label{sp:ram}
\end{align}
The bounds on the first two lines are the source's \source{New/Spanning/SecondMoment/SpanningTheorem}{257}{near-isometry} and \source{New/Spanning/SecondMoment/SpanningTheorem}{288}{second-moment estimate}. The \source{New/Reduction/ShortAndEnough}{410}{ramified-source argument} gives the third. The values are $\alpha=5b^{-4}$, $\beta=9b^{-4}$ and $\gamma=b^{-8}$.
Applying \eqref{sp:near} to $\ker(EP)$ gives
$d\ge(1-\alpha)h$, hence $h\le2d$ for $b\ge3$.

\paragraph{A single recurrence.}
Suppose the selected images span $F\subset V$, with missing fraction
$q=1-\dim F/d$. For $A_s=EU_sP|_{S_r}$ put $L_s=A_s^{-1}F$.
Adding $\operatorname{im}A_s$ gains exactly $g_s=h_r-\dim L_s$
dimensions, by rank--nullity modulo $F$. On $L_s$, the map $EU_sP$
agrees with $\Pi_F U_sP$ because its values lie in $F$ and the component
discarded by $E$ is perpendicular to $F$. Unitarity and commutation
therefore give
\[
 h_r-g_s-\alpha h
 \le \|EU_sP|_{L_s}\|_{\rm HS}^2
 \le \|\Pi_F U_sP\|_{\rm HS}^2 .
\]
This use of the actual source $S_r$ charges its missing dimensions once
in the growth law.

Let $D=\mathbb E_s U_sPP^*U_s^*$. Expanding the square and grouping
pairs of elements by $s^{-1}t$ gives
$\operatorname{tr}(D^2)=\mathbb E_s\|P^*U_sP\|_{\rm HS}^2$.
The average of the last squared norm in the preceding display is
$\operatorname{tr}(\Pi_FD)$. Trace Cauchy--Schwarz and \eqref{sp:mom}
bound this quantity by $h\sqrt{(1-q)(1+\beta)}$. Consequently
\begin{align}
 \mathbb E_s g_s
  &\ge h\{1-\alpha-\gamma-\sqrt{(1-q)(1+\beta)}\}\label{sp:exactgrowth}\\
  &\ge \tfrac h2(q-\nu),\qquad
    \nu=\beta+2\alpha+2\gamma=19b^{-4}+2b^{-8}.\nonumber
\end{align}
The last line uses $2\sqrt u\le1+u$ and $q\beta\ge0$.
Choose a translate with at least the average gain. The new fraction obeys
$q_{j+1}\le(1-h/(2d))q_j+(h/(2d))\nu$. Starting from $q_0=1$
and using $0<h/(2d)\le1$ gives, after $n$ choices,
\begin{equation}
 q_n\le\nu+(1-h/(2d))^n\le\nu+\exp(-nh/(2d)).
 \label{sp:decay}
\end{equation}
Repeated translates can be discarded without changing the span.

\paragraph{A stronger rank bound on the same budget.}
Take $n=\lceil(10d/h)\log b\rceil$ with $b\ge3$. For the resulting
set $J$ the cost and missing fraction satisfy
\begin{align}
 |J|h_r &\le nh\le 10d\log b+2d\le5d\log(1+2b^2),
   \label{sp:card}\\
 q_n&\le 19b^{-4}+b^{-5}+2b^{-8}<20b^{-4}.
   \label{sp:rank}
\end{align}
For the last inequality in \eqref{sp:card} use $5\log 2>2$.
Equations \eqref{sp:card}--\eqref{sp:rank} imply the original two
allowances, since $20b^{-4}\le4b^{-2}$ for $b\ge3$.
The coordinate isomorphism and the nonzero scalar normalizations identify
these images of $S_r$ with the images used by Theorem~B. Independent
vectors can still be selected from the actual translates of the short
integral forms supplied by Theorem~A.

\paragraph{The sharp obstruction from the two moments.}
Take $F$ to be the sum of \emph{all} allowed images. Then $g_s=0$ for
every $s$, so \eqref{sp:exactgrowth} gives a permanent missed fraction of
at most
\[
 q_* = 1-\frac{(1-\alpha-\gamma)^2}{1+\beta}
 \qquad(\alpha+\gamma<1).
\]
For full source this threshold can hold with equality. Let $b\ge3$ be
an integer, $x=b^4$, $m=(x-5)^2$, and $d=x(x+9)$. On $V=H=\mathbb C^d$,
a cyclic group of order $m$ permutes the first $m$ coordinate axes
regularly and fixes the other $d-m$. Take $E=I$, $S_r=S=\mathbb C$
and $P(1)=\sqrt{1-5/x}\,e_1$. The near-isometry condition is sharp.
Only the identity group element has $P^*U_sP\ne0$, so the normalized
second moment is
\[
 d\,\mathbb E_s\|P^*U_sP\|_{\rm HS}^2
       =\frac{d(1-5/x)^2}{m}=1+9/x.
\]
The whole orbit spans exactly the first $m$ axes. Its missed fraction is
\[
 \frac{19x-25}{x(x+9)}
       =1-\frac{(1-5/x)^2}{1+9/x}
       =(19+o(1))b^{-4}.
\]
Thus any universal leading coefficient smaller than $19$ fails for these
three input inequalities, even with access to every translate. Further
arithmetic information about the Poincar\'e map could produce a stronger
conclusion in the actual automorphic spaces. The example determines the
limit of the moment argument. For a concrete comparison of the abstract estimates, at $b=10$ the published allowance is $4\%$, the new upper bound is $0.191002\%$, and this cyclic example actually misses about $0.189804\%$.

\paragraph{Relation to the existing proof.}
The group average, Cauchy--Schwarz bound, and preimage rank identity are
already present in \texttt{FewTranslates.lean}. The change is to apply
them on $S_r$ in the first step, retain the original numerical errors,
and solve one affine recurrence. The existing proof first weakens both
moment errors to $(8b^2)^{-1}$, obtains enough translates of the entire
source, and then bounds the cumulative loss on restricting their images.
The local orbital estimates enter at precisely \eqref{sp:near} and \eqref{sp:mom}. Thus a formal refactor can reuse the source's \source{New/Spanning/SecondMoment/FewTranslates}{288}{Hilbert--Schmidt estimate} and its \source{New/Spanning/SecondMoment/FewTranslates}{393}{rank identity}, then change the dimension recurrence and the final restriction step.
